% ══════════════════════════════════════════════════════════════════════════════
%                         فصل هفتم
%                 سوسیال‌دموکراسی و چپ میانه
% ══════════════════════════════════════════════════════════════════════════════

\chapter{سوسیال‌دموکراسی و چپ میانه}
\label{ch:social-democracy}

\begin{flushright}
\textit{«هدف نهایی، هرچه باشد، برای من هیچ است؛ جنبش همه چیز است.»}

— ادوارد برنشتاین (۱۸۹۹)
\end{flushright}

% ══════════════════════════════════════════════════════════════════════════════
%                         مقدمه فصل
% ══════════════════════════════════════════════════════════════════════════════

\lettrine[lines=3, lraise=0.1, nindent=0.5em]{\textcolor{CenterGreen}{س}}{وسیال‌دموکراسی} 
شاید موفق‌ترین شاخه سوسیالیسم در تاریخ باشد—اگر موفقیت را با تأثیر بر زندگی واقعی میلیون‌ها انسان بسنجیم. دولت رفاه، بیمه بیکاری، بازنشستگی، بهداشت عمومی، آموزش رایگان—همه این‌ها عمدتاً دستاورد احزاب سوسیال‌دموکرات در قرن بیستم‌اند. اما این موفقیت با قیمتی به دست آمد: کنار گذاشتن آرمان انقلاب و پذیرش سرمایه‌داری. این فصل داستان این سازش تاریخی را روایت می‌کند.

% ══════════════════════════════════════════════════════════════════════════════
\section{سوسیال‌دموکراسی چیست؟}
% ══════════════════════════════════════════════════════════════════════════════

\subsection{تعریف و تحول تاریخی}

\begin{defbox}[تعریف: سوسیال‌دموکراسی]
واژه «سوسیال‌دموکراسی» در طول زمان معناهای متفاوتی داشته:

\textbf{معنای اولیه (قرن ۱۹):}
\begin{itemize}[nosep]
    \item مترادف با مارکسیسم
    \item SPD آلمان «حزب سوسیال‌دموکرات» بود و مارکسیست
\end{itemize}

\textbf{معنای کنونی (پس از ۱۹۱۴):}
\begin{itemize}[nosep]
    \item گرایش اصلاح‌طلب در سوسیالیسم
    \item پذیرش دموکراسی پارلمانی و اقتصاد مختلط
    \item هدف: اصلاح سرمایه‌داری، نه سرنگونی آن
    \item ابزار: دولت رفاه، مالیات تصاعدی، تنظیم بازار
\end{itemize}
\end{defbox}

\begin{figure}[htbp]
\centering
\begin{tikzpicture}[scale=0.85]

% عنوان
\node[font=\large\bfseries, text=DarkGray] at (0, 7) {\fa{تحول سوسیال‌دموکراسی از مارکسیسم به میانه}};

% محور زمان
\draw[line width=4pt, MediumGray] (-7, 4.5) -- (7, 4.5);

% نقاط
% مارکسیسم
\fill[LeftRedDark] (-6, 4.5) circle (10pt);
\node[above, text=LeftRedDark, font=\small\bfseries] at (-6, 5) {\fa{۱۸۷۵}};
\node[below, text width=2.2cm, align=center, text=DarkGray, font=\scriptsize] at (-6, 3.8) {
    \fa{تأسیس SPD}\\
    \fa{مارکسیسم ارتدکس}
};

% برنشتاین
\fill[LeftRed] (-3, 4.5) circle (10pt);
\node[above, text=LeftRed, font=\small\bfseries] at (-3, 5) {\fa{۱۸۹۹}};
\node[below, text width=2.2cm, align=center, text=DarkGray, font=\scriptsize] at (-3, 3.8) {
    \fa{رویزیونیسم}\\
    \fa{برنشتاین}
};

% انشعاب
\fill[OrangeAccent] (0, 4.5) circle (10pt);
\node[above, text=OrangeAccent, font=\small\bfseries] at (0, 5) {\fa{۱۹۱۴-۱۹۱۷}};
\node[below, text width=2.5cm, align=center, text=DarkGray, font=\scriptsize] at (0, 3.7) {
    \fa{جنگ + انقلاب روسیه}\\
    \fa{انشعاب بزرگ}
};

% دولت رفاه
\fill[CenterGreen] (3, 4.5) circle (10pt);
\node[above, text=CenterGreen, font=\small\bfseries] at (3, 5) {\fa{۱۹۴۵-۱۹۷۵}};
\node[below, text width=2.2cm, align=center, text=DarkGray, font=\scriptsize] at (3, 3.8) {
    \fa{دولت رفاه}\\
    \fa{عصر طلایی}
};

% بادگودسبرگ
\fill[CenterGreenLight] (6, 4.5) circle (10pt);
\node[above, text=CenterGreenLight, font=\small\bfseries] at (6, 5) {\fa{۱۹۵۹}};
\node[below, text width=2.2cm, align=center, text=DarkGray, font=\scriptsize] at (6, 3.8) {
    \fa{بادگودسبرگ}\\
    \fa{رسماً غیرمارکسیست}
};

% فلش حرکت
\draw[line width=2pt, color=MediumGray, -{Stealth[length=8pt]}] (-6, 2.5) -- (6, 2.5);
\node[text=LeftRed, font=\small] at (-5, 2) {\fa{چپ}};
\node[text=CenterGreen, font=\small] at (5, 2) {\fa{میانه}};

\end{tikzpicture}
\caption{سیر تحول سوسیال‌دموکراسی}
\label{fig:socdem-evolution}
\end{figure}

\subsection{ویژگی‌های کلیدی}

\begin{table}[htbp]
\centering
\caption{ویژگی‌های سوسیال‌دموکراسی معاصر}
\label{tab:socdem-features}
\begin{tabular}{r|r}
\toprule
\textbf{حوزه} & \textbf{موضع سوسیال‌دموکراتیک} \\
\midrule
اقتصاد & اقتصاد مختلط؛ بازار + تنظیم دولتی \\
مالکیت & عمدتاً خصوصی؛ برخی صنایع کلیدی دولتی \\
رفاه & دولت رفاه گسترده؛ بیمه، بهداشت، آموزش \\
مالیات & تصاعدی؛ بازتوزیع ثروت \\
کار & حمایت از اتحادیه‌ها؛ حقوق کارگری \\
روش تغییر & اصلاحات تدریجی از طریق دموکراسی \\
انقلاب & رد شده \\
\bottomrule
\end{tabular}
\end{table}

% ══════════════════════════════════════════════════════════════════════════════
\section{مناظره برنشتاین: نقطه عطف}
% ══════════════════════════════════════════════════════════════════════════════

\begin{personbox}[ادوارد برنشتاین (۱۸۵۰-۱۹۳۲)]

\textbf{ملیت:} آلمانی

\textbf{اثر کلیدی:} \book{پیش‌شرط‌های سوسیالیسم و وظایف سوسیال‌دموکراسی} (۱۸۹۹)

\textbf{ایده محوری:} رویزیونیسم (بازنگری در مارکسیسم)

\mediumspace

\person{برنشتاین}، که با مارکس و انگلس نزدیک بود، در اواخر قرن نوزدهم نقدی بنیادین به مارکسیسم ارتدکس وارد کرد:

\textbf{ادعاهای برنشتاین:}
\begin{enumerate}[nosep]
    \item \textbf{طبقه متوسط محو نشده}—برخلاف پیش‌بینی مارکس، رشد کرده
    \item \textbf{بحران‌های سرمایه‌داری} کمتر و ملایم‌تر شده‌اند
    \item \textbf{وضعیت کارگران} بهتر شده، نه بدتر
    \item \textbf{دموکراسی} امکان اصلاحات را فراهم کرده
    \item \textbf{انقلاب خشونت‌آمیز} نه لازم است نه مطلوب
\end{enumerate}

\mediumspace

\textbf{نتیجه‌گیری:} سوسیالیسم باید از طریق اصلاحات تدریجی و دموکراتیک پیش برود.

\mediumspace

\textbf{جمله معروف:} \textit{«هدف نهایی، هرچه باشد، برای من هیچ است؛ جنبش همه چیز است.»}
\end{personbox}

\begin{debatebox}[واکنش‌ها به رویزیونیسم برنشتاین]
\textbf{کارل کائوتسکی (مارکسیسم ارتدکس):}
\begin{itemize}[nosep]
    \item مارکسیسم نظریه علمی است، نباید تغییر کند
    \item انقلاب «اجتناب‌ناپذیر» است
    \item اصلاحات خوب است، اما کافی نیست
\end{itemize}

\textbf{رزا لوکزامبورگ:}
\begin{itemize}[nosep]
    \item اصلاحات در چارچوب سرمایه‌داری محدودند
    \item برنشتاین به لیبرالیسم تسلیم شده
    \item «اصلاح یا انقلاب»—باید انقلاب را انتخاب کرد
\end{itemize}

\textbf{لنین:}
\begin{itemize}[nosep]
    \item رویزیونیسم = خیانت طبقاتی
    \item «آنارشیسم کودکانه» نقش حزب را نفی می‌کند
    \item برنشتاین عامل بورژوازی در جنبش کارگری است
\end{itemize}

\textbf{اما در عمل:} SPD همین راه را رفت، هرچند رسماً انکار کرد.
\end{debatebox}

% ══════════════════════════════════════════════════════════════════════════════
\section{مدل اسکاندیناوی}
% ══════════════════════════════════════════════════════════════════════════════

موفق‌ترین تجربه سوسیال‌دموکراسی در کشورهای اسکاندیناوی—به‌ویژه سوئد—تحقق یافت.

\subsection{دولت رفاه سوئدی}

\begin{historybox}[مدل سوئدی (Folkhemmet)]
\textbf{فولک‌همت} (خانه مردم): استعاره‌ای که سوسیال‌دموکرات‌های سوئد به کار بردند.

\textbf{ویژگی‌ها:}
\begin{itemize}[nosep]
    \item دولت رفاه جامع: از گهواره تا گور
    \item مالیات بالا (گاهی بیش از ۵۰٪)
    \item خدمات عمومی با کیفیت: بهداشت، آموزش، مراقبت از کودک
    \item بازار کار منعطف با حمایت قوی از بیکاران
    \item \textbf{کورپوراتیسم:} مذاکره سه‌جانبه دولت، کارفرما، اتحادیه
\end{itemize}

\textbf{SAP (حزب سوسیال‌دموکرات):} از ۱۹۳۲ تا ۱۹۷۶ بدون وقفه حکومت کرد.

\textbf{نتایج:} یکی از برابرترین و مرفه‌ترین جوامع جهان
\end{historybox}

\begin{figure}[htbp]
\centering
\begin{tikzpicture}[scale=0.85]

% عنوان
\node[font=\large\bfseries, text=DarkGray] at (0, 7) {\fa{ساختار دولت رفاه اسکاندیناوی}};

% مرکز
\fill[CenterGreen, opacity=0.9] (0, 3.5) circle (2cm);
\node[text=white, font=\bfseries\large] at (0, 3.5) {\fa{دولت رفاه}};

% شاخه‌ها
% بهداشت
\fill[CenterGreenLight, opacity=0.8] (-4, 5.5) circle (1.2cm);
\node[text=DarkGray, font=\small\bfseries] at (-4, 5.5) {\fa{بهداشت}};
\draw[line width=2pt, color=CenterGreen] (-2, 4.3) -- (-3, 4.8);

% آموزش
\fill[CenterGreenLight, opacity=0.8] (4, 5.5) circle (1.2cm);
\node[text=DarkGray, font=\small\bfseries] at (4, 5.5) {\fa{آموزش}};
\draw[line width=2pt, color=CenterGreen] (2, 4.3) -- (3, 4.8);

% بازنشستگی
\fill[CenterGreenLight, opacity=0.8] (-4.5, 2) circle (1.2cm);
\node[text=DarkGray, font=\small\bfseries] at (-4.5, 2) {\fa{بازنشستگی}};
\draw[line width=2pt, color=CenterGreen] (-2, 3) -- (-3.3, 2.5);

% بیمه بیکاری
\fill[CenterGreenLight, opacity=0.8] (4.5, 2) circle (1.2cm);
\node[text=DarkGray, font=\small\bfseries, text width=1.8cm, align=center] at (4.5, 2) {\fa{بیمه}\\[-2pt]\fa{بیکاری}};
\draw[line width=2pt, color=CenterGreen] (2, 3) -- (3.3, 2.5);

% مراقبت کودک
\fill[CenterGreenLight, opacity=0.8] (0, 0.5) circle (1.2cm);
\node[text=DarkGray, font=\small\bfseries, text width=1.8cm, align=center] at (0, 0.5) {\fa{مراقبت}\\[-2pt]\fa{کودک}};
\draw[line width=2pt, color=CenterGreen] (0, 1.5) -- (0, 1.7);

% تأمین مالی
\node[text=MediumGray, font=\small, text width=8cm, align=center] at (0, -1) {
    \fa{تأمین مالی: مالیات تصاعدی بالا + مشارکت کارفرما و کارگر}
};

\end{tikzpicture}
\caption{اجزای دولت رفاه اسکاندیناوی}
\label{fig:scandinavian-welfare}
\end{figure}

\begin{personbox}[اولوف پالمه (۱۹۲۷-۱۹۸۶)]

\textbf{ملیت:} سوئدی

\textbf{سمت:} نخست‌وزیر سوئد (۱۹۶۹-۱۹۷۶، ۱۹۸۲-۱۹۸۶)

\textbf{ایده محوری:} سوسیال‌دموکراسی رادیکال و بین‌الملل‌گرایی

\mediumspace

\person{پالمه} چهره برجسته سوسیال‌دموکراسی جهانی بود:

\begin{itemize}[nosep]
    \item گسترش دولت رفاه
    \item نقد جنگ ویتنام و امپریالیسم آمریکا
    \item حمایت از جنبش‌های رهایی‌بخش جهان سوم
    \item میانجی‌گری در جنگ ایران و عراق
    \item دفاع از خلع سلاح هسته‌ای
\end{itemize}

\mediumspace

\textbf{سرنوشت:} ترور در خیابان استکهلم (۱۹۸۶)؛ قاتل هنوز رسماً شناسایی نشده.

\mediumspace

\textbf{نقل‌قول:} \textit{«سیاست یعنی خواستن.»}
\end{personbox}

\subsection{کورپوراتیسم اجتماعی}

\begin{defbox}[تعریف: کورپوراتیسم اجتماعی (نئوکورپوراتیسم)]
نظامی که در آن:
\begin{itemize}[nosep]
    \item \textbf{دولت، اتحادیه‌های کارگری و انجمن‌های کارفرمایی} مذاکره می‌کنند
    \item تصمیمات اقتصادی کلان با توافق گرفته می‌شود
    \item اتحادیه‌ها در ازای دستمزد معقول، صلح صنعتی تضمین می‌کنند
    \item کارفرمایان سرمایه‌گذاری و اشتغال را تضمین می‌کنند
    \item دولت رفاه و ثبات را تضمین می‌کند
\end{itemize}

\textbf{تفاوت با کورپوراتیسم فاشیستی:} داوطلبانه، دموکراتیک، حفظ استقلال اتحادیه‌ها
\end{defbox}

% ══════════════════════════════════════════════════════════════════════════════
\section{حزب کارگر بریتانیا}
% ══════════════════════════════════════════════════════════════════════════════

\begin{figure}[htbp]
\centering
\begin{tikzpicture}[scale=0.85]

% عنوان
\node[font=\large\bfseries, text=DarkGray] at (0, 7.5) {\fa{تایم‌لاین: حزب کارگر بریتانیا}};

% خط زمان
\draw[line width=4pt, MediumGray] (-7, 5.5) -- (7, 5.5);

% رویدادها
% تأسیس
\fill[LeftRedLight] (-6, 5.5) circle (9pt);
\node[above, text=LeftRedLight, font=\scriptsize\bfseries] at (-6, 6) {\fa{۱۹۰۰}};
\node[below, text width=2cm, align=center, text=DarkGray, font=\tiny] at (-6, 4.9) {
    \fa{تأسیس}\\
    \fa{(اتحادیه‌ها+فابیان)}
};

% اولین دولت
\fill[LeftRed] (-4, 5.5) circle (9pt);
\node[above, text=LeftRed, font=\scriptsize\bfseries] at (-4, 6) {\fa{۱۹۲۴}};
\node[below, text width=2cm, align=center, text=DarkGray, font=\tiny] at (-4, 4.9) {
    \fa{اولین دولت}\\
    \fa{(مک‌دونالد)}
};

% اتلی
\fill[CenterGreen] (-1.5, 5.5) circle (11pt);
\node[above, text=CenterGreen, font=\scriptsize\bfseries] at (-1.5, 6.1) {\fa{۱۹۴۵-۵۱}};
\node[below, text width=2.2cm, align=center, text=DarkGray, font=\tiny] at (-1.5, 4.8) {
    \fa{دولت اتلی}\\
    \fa{NHS، ملی‌سازی}
};

% ویلسون
\fill[CenterGreenLight] (1, 5.5) circle (9pt);
\node[above, text=CenterGreenLight, font=\scriptsize\bfseries] at (1, 6) {\fa{۶۴-۷۰، ۷۴-۷۶}};
\node[below, text width=2cm, align=center, text=DarkGray, font=\tiny] at (1, 4.9) {
    \fa{ویلسون}\\
    \fa{اصلاحات اجتماعی}
};

% بحران
\fill[OrangeAccent] (3.5, 5.5) circle (9pt);
\node[above, text=OrangeAccent, font=\scriptsize\bfseries] at (3.5, 6) {\fa{۱۹۷۹-۹۷}};
\node[below, text width=2cm, align=center, text=DarkGray, font=\tiny] at (3.5, 4.9) {
    \fa{اپوزیسیون}\\
    \fa{(عصر تاچر)}
};

% بلر
\fill[RightBlueLight] (6, 5.5) circle (10pt);
\node[above, text=RightBlueLight, font=\scriptsize\bfseries] at (6, 6) {\fa{۱۹۹۷-۲۰۱۰}};
\node[below, text width=2cm, align=center, text=DarkGray, font=\tiny] at (6, 4.8) {
    \fa{نیو لیبر}\\
    \fa{(بلر، براون)}
};

\end{tikzpicture}
\caption{سیر تاریخی حزب کارگر بریتانیا}
\label{fig:labour-timeline}
\end{figure}

\subsection{دولت اتلی و بنیان‌گذاری دولت رفاه}

\begin{personbox}[کلمنت اتلی (۱۸۸۳-۱۹۶۷)]

\textbf{ملیت:} بریتانیایی

\textbf{سمت:} نخست‌وزیر (۱۹۴۵-۱۹۵۱)

\textbf{ایده محوری:} سوسیالیسم دموکراتیک عملی

\mediumspace

\person{اتلی} در انتخابات ۱۹۴۵، چرچیل قهرمان جنگ را شکست داد و دولتی تشکیل داد که بریتانیا را دگرگون کرد:

\textbf{دستاوردها:}
\begin{itemize}[nosep]
    \item \textbf{NHS:} خدمات ملی بهداشت—بهداشت رایگان برای همه
    \item \textbf{ملی‌سازی:} ذغال‌سنگ، راه‌آهن، فولاد، برق، گاز
    \item \textbf{دولت رفاه:} بر اساس گزارش بوریج
    \item \textbf{آموزش:} گسترش آموزش رایگان
    \item \textbf{مسکن:} ساخت میلیون‌ها خانه دولتی
\end{itemize}

\mediumspace

\textbf{نقل‌قول:} \textit{«دموکراسی یعنی حکومت از طریق بحث، اما فقط وقتی کارآمد است که بتوانی مردم را وادار به سکوت کنی.»}
\end{personbox}

\begin{historybox}[NHS: خدمات ملی بهداشت (۱۹۴۸)]
\textbf{بنیان‌گذار:} آنورین بوان (وزیر بهداشت)

\textbf{اصول:}
\begin{itemize}[nosep]
    \item بهداشت رایگان در نقطه استفاده
    \item تأمین مالی از مالیات عمومی
    \item دسترسی همگانی بدون توجه به توانایی پرداخت
\end{itemize}

\textbf{مقاومت:} انجمن پزشکان مخالف بود. بوان گفت: «دهان آن‌ها را با طلا پر کردم.»

\textbf{میراث:} NHS هنوز محبوب‌ترین نهاد بریتانیاست. حتی تاچر جرأت نکرد آن را خصوصی کند.
\end{historybox}

\subsection{از ویلسون تا بلر}

\begin{table}[htbp]
\centering
\caption{تحول حزب کارگر: از چپ به میانه}
\label{tab:labour-evolution}
\begin{tabular}{r|c|c|c}
\toprule
\textbf{دوره} & \textbf{رهبر} & \textbf{موضع} & \textbf{سیاست کلیدی} \\
\midrule
۱۹۴۵-۵۱ & اتلی & چپ سنتی & ملی‌سازی، NHS \\
۱۹۶۴-۷۰ & ویلسون & چپ میانه & اصلاحات اجتماعی \\
۱۹۷۴-۷۹ & ویلسون/کالاهان & بحران & بحران اقتصادی \\
۱۹۸۳ & فوت & چپ رادیکال & شکست سنگین \\
۱۹۹۷-۲۰۱۰ & بلر/براون & راه سوم & بازار + عدالت اجتماعی \\
۲۰۱۵-۲۰ & کوربین & چپ & بازگشت به ریشه‌ها \\
\bottomrule
\end{tabular}
\end{table}

% ══════════════════════════════════════════════════════════════════════════════
\section{سوسیال‌دموکراسی آلمان}
% ══════════════════════════════════════════════════════════════════════════════

SPD (حزب سوسیال‌دموکرات آلمان) قدیمی‌ترین و تأثیرگذارترین حزب سوسیالیست تاریخ است.

\begin{figure}[htbp]
\centering
\begin{tikzpicture}[scale=0.85]

% عنوان
\node[font=\large\bfseries, text=DarkGray] at (0, 7.5) {\fa{تایم‌لاین: تحول SPD آلمان}};

% خط زمان
\draw[line width=4pt, MediumGray] (-7, 5.5) -- (7, 5.5);

% رویدادها
% تأسیس
\fill[LeftRedDark] (-6, 5.5) circle (9pt);
\node[above, text=LeftRedDark, font=\scriptsize\bfseries] at (-6, 6) {\fa{۱۸۷۵}};
\node[below, text width=2cm, align=center, text=DarkGray, font=\tiny] at (-6, 4.9) {
    \fa{تأسیس}\\
    \fa{گوتا}
};

% ارفورت
\fill[LeftRed] (-4, 5.5) circle (9pt);
\node[above, text=LeftRed, font=\scriptsize\bfseries] at (-4, 6) {\fa{۱۸۹۱}};
\node[below, text width=2cm, align=center, text=DarkGray, font=\tiny] at (-4, 4.9) {
    \fa{برنامه ارفورت}\\
    \fa{مارکسیسم}
};

% جنگ
\fill[OrangeAccent] (-2, 5.5) circle (9pt);
\node[above, text=OrangeAccent, font=\scriptsize\bfseries] at (-2, 6) {\fa{۱۹۱۴}};
\node[below, text width=2cm, align=center, text=DarkGray, font=\tiny] at (-2, 4.9) {
    \fa{رأی به جنگ}\\
    \fa{انشعاب}
};

% وایمار
\fill[CenterGreenLight] (0, 5.5) circle (9pt);
\node[above, text=CenterGreenLight, font=\scriptsize\bfseries] at (0, 6) {\fa{۱۹۱۸-۳۳}};
\node[below, text width=2cm, align=center, text=DarkGray, font=\tiny] at (0, 4.9) {
    \fa{جمهوری وایمار}\\
    \fa{دولت‌داری}
};

% ممنوعیت
\fill[MediumGray] (2, 5.5) circle (9pt);
\node[above, text=MediumGray, font=\scriptsize\bfseries] at (2, 6) {\fa{۱۹۳۳-۴۵}};
\node[below, text width=2cm, align=center, text=DarkGray, font=\tiny] at (2, 4.9) {
    \fa{نازیسم}\\
    \fa{ممنوعیت}
};

% بادگودسبرگ
\fill[CenterGreen] (4.5, 5.5) circle (11pt);
\node[above, text=CenterGreen, font=\scriptsize\bfseries] at (4.5, 6.1) {\fa{۱۹۵۹}};
\node[below, text width=2.3cm, align=center, text=DarkGray, font=\tiny] at (4.5, 4.8) {
    \fa{بادگودسبرگ}\\
    \fa{غیرمارکسیست}
};

% برانت
\fill[CenterGreenLight] (6.5, 5.5) circle (9pt);
\node[above, text=CenterGreenLight, font=\scriptsize\bfseries] at (6.5, 6) {\fa{۱۹۶۹-۸۲}};
\node[below, text width=2cm, align=center, text=DarkGray, font=\tiny] at (6.5, 4.9) {
    \fa{برانت/اشمیت}\\
    \fa{قدرت}
};

\end{tikzpicture}
\caption{سیر تاریخی حزب سوسیال‌دموکرات آلمان}
\label{fig:spd-timeline}
\end{figure}

\subsection{برنامه بادگودسبرگ (۱۹۵۹)}

\begin{historybox}[برنامه بادگودسبرگ: نقطه عطف]
در ۱۹۵۹، SPD در کنگره بادگودسبرگ برنامه جدیدی تصویب کرد که رسماً با مارکسیسم خداحافظی می‌کرد:

\textbf{کنار گذاشته شد:}
\begin{itemize}[nosep]
    \item مارکسیسم به عنوان ایدئولوژی رسمی
    \item ملی‌سازی گسترده
    \item مبارزه طبقاتی به عنوان موتور تاریخ
\end{itemize}

\textbf{پذیرفته شد:}
\begin{itemize}[nosep]
    \item اقتصاد بازار اجتماعی
    \item مالکیت خصوصی (با مسئولیت اجتماعی)
    \item دفاع ملی و عضویت در ناتو
    \item ارزش‌های مسیحی، انسان‌گرایی و سوسیالیسم دموکراتیک
\end{itemize}

\textbf{شعار:} «رقابت تا حد ممکن، برنامه‌ریزی تا حد لازم»

\textbf{نتیجه:} SPD از ۱۹۶۶ وارد دولت شد و از ۱۹۶۹ تا ۱۹۸۲ صدراعظم داشت.
\end{historybox}

\begin{personbox}[ویلی برانت (۱۹۱۳-۱۹۹۲)]

\textbf{ملیت:} آلمانی

\textbf{سمت:} صدراعظم آلمان غربی (۱۹۶۹-۱۹۷۴)

\textbf{جایزه صلح نوبل:} ۱۹۷۱

\mediumspace

\person{برانت} نمادی از سوسیال‌دموکراسی آلمانی پس از جنگ بود:

\begin{itemize}[nosep]
    \item \textbf{اوست‌پولیتیک:} سیاست نزدیکی با شرق
    \item زانو زدن در وارشاو در برابر بنای یادبود قیام گتو (۱۹۷۰)
    \item اصلاحات داخلی: آموزش، حقوق زنان
    \item رئیس انترناسیونال سوسیالیست (۱۹۷۶-۱۹۹۲)
\end{itemize}

\mediumspace

\textbf{نقل‌قول:} \textit{«سیاست صلح، سیاست شناختن واقعیت است.»}
\end{personbox}

% ══════════════════════════════════════════════════════════════════════════════
\section{احزاب سوسیالیست در اروپای جنوبی}
% ══════════════════════════════════════════════════════════════════════════════

\subsection{PSOE اسپانیا}

\begin{historybox}[حزب سوسیالیست کارگران اسپانیا (PSOE)]
\begin{itemize}[nosep]
    \item تأسیس: ۱۸۷۹
    \item دوران فرانکو: تبعید و زیرزمینی
    \item \person{فلیپه گونزالس}: نخست‌وزیر ۱۹۸۲-۱۹۹۶
    \item مدرن‌سازی اسپانیا، عضویت در اتحادیه اروپا
    \item از مارکسیسم به سوسیال‌دموکراسی لیبرال
\end{itemize}
\end{historybox}

\subsection{PS فرانسه و میتران}

\begin{personbox}[فرانسوا میتران (۱۹۱۶-۱۹۹۶)]

\textbf{ملیت:} فرانسوی

\textbf{سمت:} رئیس‌جمهور فرانسه (۱۹۸۱-۱۹۹۵)

\textbf{ایده محوری:} «گسست» با سرمایه‌داری

\mediumspace

\person{میتران} نخستین رئیس‌جمهور سوسیالیست جمهوری پنجم فرانسه بود:

\textbf{برنامه ۱۹۸۱ (رادیکال):}
\begin{itemize}[nosep]
    \item ملی‌سازی بانک‌ها و صنایع بزرگ
    \item افزایش حداقل دستمزد و مرخصی
    \item کاهش ساعت کار
    \item استخدام در بخش دولتی
\end{itemize}

\textbf{چرخش ۱۹۸۳:}
\begin{itemize}[nosep]
    \item فشار بازار و فرار سرمایه
    \item سیاست ریاضت اقتصادی
    \item کنار گذاشتن برنامه رادیکال
\end{itemize}

\mediumspace

\textbf{درس:} سوسیال‌دموکراسی در یک کشور، در اقتصاد جهانی، محدود است.
\end{personbox}

\subsection{PCI/PD ایتالیا}

\begin{historybox}[از کمونیسم به سوسیال‌دموکراسی]
\textbf{PCI (حزب کمونیست ایتالیا):}
\begin{itemize}[nosep]
    \item بزرگ‌ترین حزب کمونیست غرب
    \item یوروکمونیسم: فاصله از مسکو
    \item انریکو برلینگوئر: «سازش تاریخی»
\end{itemize}

\textbf{تحول:}
\begin{itemize}[nosep]
    \item ۱۹۹۱: انحلال PCI پس از فروپاشی شوروی
    \item تأسیس PDS (دموکرات‌های چپ)
    \item سپس DS و نهایتاً PD (حزب دموکراتیک)
    \item امروز: حزب میانه‌رو، نه چپ سنتی
\end{itemize}
\end{historybox}

% ══════════════════════════════════════════════════════════════════════════════
\section{بحران سوسیال‌دموکراسی}
% ══════════════════════════════════════════════════════════════════════════════

از دهه ۱۹۸۰، سوسیال‌دموکراسی با چالش‌های جدی مواجه شده است.

\begin{figure}[htbp]
\centering
\begin{tikzpicture}[scale=0.85]

% عنوان
\node[font=\large\bfseries, text=DarkGray] at (0, 7) {\fa{چالش‌های سوسیال‌دموکراسی}};

% مرکز
\fill[CenterGreen, opacity=0.8] (0, 3.5) circle (1.5cm);
\node[text=white, font=\bfseries] at (0, 3.8) {\fa{سوسیال‌}};
\node[text=white, font=\bfseries] at (0, 3.2) {\fa{دموکراسی}};

% چالش‌ها (اطراف)
% نئولیبرالیسم
\fill[RightBlue, opacity=0.7] (-4, 5.5) circle (1cm);
\node[text=white, font=\scriptsize\bfseries, text width=1.8cm, align=center] at (-4, 5.5) {\fa{نئولیبرالیسم}};
\draw[line width=2pt, color=RightBlue, -{Stealth[length=6pt]}] (-3, 4.9) -- (-1.3, 4.2);

% جهانی‌شدن
\fill[RightBlueLight, opacity=0.7] (4, 5.5) circle (1cm);
\node[text=white, font=\scriptsize\bfseries, text width=1.8cm, align=center] at (4, 5.5) {\fa{جهانی‌شدن}};
\draw[line width=2pt, color=RightBlueLight, -{Stealth[length=6pt]}] (3, 4.9) -- (1.3, 4.2);

% تغییر طبقاتی
\fill[OrangeAccent, opacity=0.7] (-4.5, 2) circle (1cm);
\node[text=white, font=\scriptsize\bfseries, text width=1.8cm, align=center] at (-4.5, 2) {\fa{زوال طبقه}\\[-2pt]\fa{کارگر}};
\draw[line width=2pt, color=OrangeAccent, -{Stealth[length=6pt]}] (-3.5, 2.5) -- (-1.4, 3.1);

% پوپولیسم
\fill[LeftRed, opacity=0.7] (4.5, 2) circle (1cm);
\node[text=white, font=\scriptsize\bfseries, text width=1.8cm, align=center] at (4.5, 2) {\fa{پوپولیسم}\\[-2pt]\fa{راست}};
\draw[line width=2pt, color=LeftRed, -{Stealth[length=6pt]}] (3.5, 2.5) -- (1.4, 3.1);

% سبزها
\fill[CenterGreenDark, opacity=0.7] (0, 0.5) circle (1cm);
\node[text=white, font=\scriptsize\bfseries, text width=1.8cm, align=center] at (0, 0.5) {\fa{رقابت}\\[-2pt]\fa{سبزها}};
\draw[line width=2pt, color=CenterGreenDark, -{Stealth[length=6pt]}] (0, 1.5) -- (0, 2);

\end{tikzpicture}
\caption{چالش‌های پیش روی سوسیال‌دموکراسی}
\label{fig:socdem-challenges}
\end{figure}

\subsection{علل بحران}

\begin{debatebox}[چرا سوسیال‌دموکراسی در بحران است؟]
\textbf{عوامل اقتصادی:}
\begin{itemize}[nosep]
    \item \textbf{جهانی‌شدن:} سرمایه می‌تواند فرار کند؛ مالیات بالا دشوار شده
    \item \textbf{نئولیبرالیسم:} از ۱۹۸۰ ایده‌های راست غالب شد
    \item \textbf{بدهی دولت:} گسترش رفاه پرهزینه شده
\end{itemize}

\textbf{عوامل اجتماعی:}
\begin{itemize}[nosep]
    \item \textbf{زوال طبقه کارگر صنعتی:} پایگاه سنتی کوچک شده
    \item \textbf{فردگرایی:} همبستگی طبقاتی کمتر شده
    \item \textbf{تنوع هویتی:} طبقه دیگر محور اصلی نیست
\end{itemize}

\textbf{عوامل سیاسی:}
\begin{itemize}[nosep]
    \item \textbf{پوپولیسم راست:} کارگران سفیدپوست به سمت راست رفته‌اند
    \item \textbf{رقابت چپ:} سبزها، چپ رادیکال
    \item \textbf{بی‌تفاوتی:} احزاب شبیه هم شده‌اند
\end{itemize}
\end{debatebox}

\begin{table}[htbp]
\centering
\caption{افت آرای احزاب سوسیال‌دموکرات (مقایسه اوج و امروز)}
\label{tab:socdem-decline}
\begin{tabular}{r|c|c|c}
\toprule
\textbf{کشور} & \textbf{حزب} & \textbf{اوج} & \textbf{اخیر} \\
\midrule
سوئد & SAP & ۵۰٪ (۱۹۶۸) & ~۳۰٪ \\
آلمان & SPD & ۴۶٪ (۱۹۷۲) & ~۲۰-۲۵٪ \\
فرانسه & PS & ۳۷٪ (۱۹۸۱) & ~۵-۷٪ (۲۰۲۲) \\
هلند & PvdA & ۳۳٪ (۱۹۷۷) & ~۶٪ \\
یونان & PASOK & ۴۸٪ (۱۹۸۱) & ~۸٪ \\
\bottomrule
\end{tabular}
\end{table}

% ══════════════════════════════════════════════════════════════════════════════
\section{اختلافات درونی}
% ══════════════════════════════════════════════════════════════════════════════

\begin{debatebox}[چپ سنتی در برابر نوسازها]
\textbf{چپ سنتی (جرمی کوربین، برنی سندرز):}
\begin{itemize}[nosep]
    \item بازگشت به ریشه‌ها: ملی‌سازی، حقوق کارگری
    \item نقد نئولیبرالیسم و جهانی‌شدن
    \item مالیات بالا بر ثروتمندان
    \item بهداشت و آموزش رایگان
    \item «مشکل این نیست که بیش از حد چپ بودیم، بلکه کم بودیم»
\end{itemize}

\textbf{نوسازها/راه سوم (بلر، شرودر):}
\begin{itemize}[nosep]
    \item پذیرش بازار و جهانی‌شدن
    \item «اقتصاد کارآمد + عدالت اجتماعی»
    \item سرمایه‌گذاری در آموزش، نه بازتوزیع مستقیم
    \item جلب طبقه متوسط
    \item «باید با زمان پیش رفت»
\end{itemize}

\textbf{نتیجه:} هر دو استراتژی تا کنون با مشکل مواجه شده‌اند.
\end{debatebox}

\begin{debatebox}[اقتصادگرایی در برابر فرهنگ‌گرایی]
\textbf{تمرکز بر اقتصاد:}
\begin{itemize}[nosep]
    \item مسائل طبقاتی و اقتصادی اولویت دارند
    \item مسائل هویتی (نژاد، جنسیت) ثانویه‌اند
    \item «ابتدا نان، بعد گل»
    \item نقد: این سیاست هویتی سفید است
\end{itemize}

\textbf{تمرکز بر فرهنگ و هویت:}
\begin{itemize}[nosep]
    \item ستم‌های غیراقتصادی (نژادپرستی، جنسیت‌گرایی) مستقل‌اند
    \item چپ باید از حقوق اقلیت‌ها دفاع کند
    \item «تقاطع‌گرایی»: ستم‌ها به هم پیوسته‌اند
    \item نقد: کارگران سفید احساس رهاشدگی می‌کنند
\end{itemize}

\textbf{این مناقشه، چپ معاصر را دوپاره کرده است.}
\end{debatebox}

% ══════════════════════════════════════════════════════════════════════════════
\section{جدول جامع: احزاب سوسیال‌دموکرات اروپا}
% ══════════════════════════════════════════════════════════════════════════════

% FIXED: Converted to portrait tabularx for better fit
\begin{table}[htbp]
\centering
\caption{احزاب سوسیال‌دموکرات اصلی اروپا}
\label{tab:socdem-parties}
\footnotesize
\begin{tabularx}{\linewidth}{r|c|c|>{\centering\arraybackslash}X|c|c}
\toprule
\textbf{کشور} & \textbf{حزب} & \textbf{تأسیس} & \textbf{دولت‌داری مهم} & \textbf{چهره} & \textbf{وضعیت} \\
\midrule
سوئد & SAP & ۱۸۸۹ & ۱۹۳۲-۱۹۷۶، ۱۹۸۲-۲۰۰۶ & پالمه & حکومت \\
\midrule
آلمان & SPD & ۱۸۷۵ & ۱۹۶۹-۸۲، ۱۹۹۸-۲۰۰۵ & برانت & حکومت \\
\midrule
بریتانیا & Labour & ۱۹۰۰ & ۴۵-۵۱، ۶۴-۷۰، ۹۷-۲۰۱۰ & اتلی & اپوزیسیون \\
\midrule
فرانسه & PS & ۱۹۶۹ & ۱۹۸۱-۱۹۹۵، ۲۰۱۲-۲۰۱۷ & میتران & بحران \\
\midrule
اسپانیا & PSOE & ۱۸۷۹ & ۱۹۸۲-۱۹۹۶، ۲۰۰۴-۲۰۱۱ & گونزالس & حکومت \\
\midrule
اتریش & SPÖ & ۱۸۸۹ & ۱۹۷۰-۲۰۰۰ & کرایسکی & اپوزیسیون \\
\bottomrule
\end{tabularx}
\end{table}

% ══════════════════════════════════════════════════════════════════════════════
%                         خلاصه و پرسش‌ها
% ══════════════════════════════════════════════════════════════════════════════

\begin{summarybox}

\textbf{۱. تعریف:} سوسیال‌دموکراسی گرایشی است که از طریق اصلاحات تدریجی و دموکراتیک، سرمایه‌داری را اصلاح می‌کند: دولت رفاه، مالیات تصاعدی، حقوق کارگری.

\textbf{۲. مناظره برنشتاین:} رویزیونیسم استدلال کرد انقلاب نه لازم است نه مطلوب. در عمل، سوسیال‌دموکراسی همین راه را رفت.

\textbf{۳. مدل اسکاندیناوی:} دولت رفاه جامع، مالیات بالا، کورپوراتیسم. سوئد موفق‌ترین نمونه.

\textbf{۴. بریتانیا:} دولت اتلی NHS و ملی‌سازی را ایجاد کرد. بلر با «نیو لیبر» به راست چرخید.

\textbf{۵. آلمان:} SPD با برنامه بادگودسبرگ (۱۹۵۹) رسماً غیرمارکسیست شد. برانت و اوست‌پولیتیک.

\textbf{۶. اروپای جنوبی:} PSOE، PS فرانسه، PCI ایتالیا همگی مسیر مشابهی طی کردند: از رادیکالیسم به میانه.

\textbf{۷. بحران:} جهانی‌شدن، نئولیبرالیسم، زوال طبقه کارگر، پوپولیسم راست همگی سوسیال‌دموکراسی را به چالش کشیده‌اند.

\textbf{۸. اختلافات درونی:} چپ سنتی vs نوسازها؛ اقتصادگرایی vs فرهنگ‌گرایی.

\end{summarybox}

\vspace{5pt}

\begin{questionbox}

\begin{enumerate}[nosep, rightmargin=1cm]
    \item آیا برنشتاین درباره ضرورت نداشتن انقلاب حق داشت؟ یا سوسیال‌دموکراسی فقط «سرمایه‌داری با چهره انسانی» است؟

    \item مدل اسکاندیناوی چقدر قابل تعمیم است؟ آیا کشورهای بزرگ‌تر یا متنوع‌تر می‌توانند همین مسیر را طی کنند؟

    \item آیا بحران سوسیال‌دموکراسی ناشی از «بیش از حد به راست رفتن» است یا «کافی نبودن تطبیق با واقعیت‌های جدید»؟

    \item در عصر جهانی‌شدن، آیا سوسیال‌دموکراسی در یک کشور اصلاً ممکن است؟
\end{enumerate}

\end{questionbox}

% ══════════════════════════════════════════════════════════════════════════════
%                         پایان فصل هفتم
% ══════════════════════════════════════════════════════════════════════════════